\chapter{Random Oracle Programming: Paper Proof and Code Analysis}
\label{ch:rom-programming}

\section{Lazy ROM State}

The random oracle is modeled by a finite map and query log.  On a new query it
samples and records an answer; on a repeated query it returns the stored answer.
This lazy representation allows the proof to state exactly what the adversary
has queried and exactly when programming is safe.

Formally, the lazy ROM invariant says that the table \(T\) is a partial
realization of an ideal function.  If \(x\notin\operatorname{dom}(T)\), then the
conditional law of \(H(x)\), given the current transcript, is still the declared
output law \(D_x\).  If \(x\in\operatorname{dom}(T)\), then consistency requires
all future calls on \(x\) to return the stored value.  Thus the table domain is
the mathematical object that separates fresh randomness from already observed
randomness.

For HAETAE, the input space is not a single untyped string space in the proof.
The EasyCrypt model uses typed constructors for message hashing, challenge
hashing, matrix expansion, and sampler expansion.  This typed interface is a
formal domain-separation discipline.  It prevents a proof about one oracle site
from accidentally applying to another site with a different mathematical role.

\section{Pen-and-Paper Programming Lemma}

\begin{lemma}[ROM programming at a fresh point]
Let \(H\) be a random oracle and let \(x\) be a point sampled independently of
an adversary's prior query set \(Q\).  If \(|Q|\le q_H\) and \(x\) has support
min-entropy at least \(\log_2 N\) against the adversary's view, equivalently
every point has conditional probability at most \(1/N\), then
\[
\Pr[x\in Q] \le \frac{q_H}{N}.
\]
Conditioned on \(x\notin Q\), assigning a value to \(H(x)\) changes no
adversary-visible answer to any previous query.  If the assigned value is sampled
from the same conditional output law \(D_x\), independently of the prior view,
then the future lazy-oracle interface has the same distribution as ordinary lazy
sampling at \(x\).  If instead the simulator assigns a fixed or otherwise
biased value, freshness alone gives only the previous-answer preservation fact;
the distributional effect of future queries to \(x\) must be justified by the
reduction or accounted for as a separate game-hop condition.
\end{lemma}

\begin{proof}
For any fixed query set \(Q\), the union bound gives
\(\Pr[x\in Q]\le \sum_{u\in Q}\Pr[x=u]\).  If every point has probability at
most \(1/N\), this is at most \(|Q|/N\le q_H/N\).  If \(x\notin Q\), the lazy
random oracle has not assigned a value to \(x\).  Assigning a value at \(x\)
therefore leaves the map entries for all previous queries unchanged.  When the
assigned value has law \(D_x\), the programmed table is distributed exactly as if
the lazy oracle had sampled \(H(x)\) on its first future query.  Without that
same-law condition, the lemma should be used only as a preservation statement for
the already observed oracle answers.
\end{proof}

\section{Programming Sites}

A programming site is the oracle input where the simulator wants to control the
challenge.  The file \code{HAETAE\_ROM\_Programming.ec} proves that honest
signing produces programming sites with the necessary structural and freshness
properties.  Lemmas such as
\code{programming\_site\_self\_no\_reprogramming\_failure} and
\code{honest\_signing\_programming\_site\_query\_spread\_bound} capture the
core obligations.

The relevant mathematical object is a random variable \(X_{\mathsf{prog}}\)
whose value is the site to be programmed.  The proof must establish three facts:
the site equation identifying \(X_{\mathsf{prog}}\) with the query produced by
the signing or challenge construction, the freshness or spread bound against the
adversary's log, and the preservation statement saying that programming
\(X_{\mathsf{prog}}\) changes no previously visible oracle answer.  Missing any
one of these facts invalidates the programming step.

\section{Detailed EasyCrypt Code Analysis}

\code{HAETAE\_ROM\_Programming.ec} separates four concerns that are often merged
in paper proofs.

First, the counted lazy ROM lemmas, such as
\code{counted\_lazy\_rom\_init\_clears} and
\code{counted\_lazy\_rom\_get\_preserves\_clear}, establish basic state
invariants.  These lemmas say that initialization clears logs and that ordinary
queries do not accidentally set programming-failure flags.

Second, programming-site lemmas characterize the exact query that honest
signing will need.  For example,
\code{honest\_signing\_programming\_site\_queryE} is an equation-style lemma:
it identifies the formal query produced by the signing algorithm.  Without such
lemmas, later freshness statements would refer to a site whose connection to
signing is only informal.

Third, spread and budget lemmas translate entropy into probability bounds.
\code{honest\_signing\_programming\_site\_query\_spread\_bound} is the
machine-checked counterpart of the paper inequality \(\Pr[x\in Q]\le q_H/N\).
The EasyCrypt proof must explicitly name the distribution, support, and query
log.

Fourth, bad-event lemmas connect the local ROM facts to game-level losses.
Lemmas such as
\code{programmed\_site\_prequery\_one\_step\_challenge\_concrete\_bound} and
\code{programmed\_site\_prequery\_one\_step\_message\_hash\_counter\_bound}
are the concrete interface consumed by the hop games.

\section{Bad Events and Fundamental Lemma Use}

The proof makes bad events explicit.  Prequery events, reprogramming conflicts,
and entropy failures are represented by predicates or flags.  Fundamental-lemma
arguments then bound the difference between games by the probability of these
bad events.  In the EasyCrypt code, this means proving both a relational
similarity condition before the bad event and a probability bound for setting
that event.

This is the point where informal proofs most often hide work.  A statement such
as ``the adversary did not query this point except with probability \(q_H/N\)''
contains an independence claim, a support-size claim, a query-budget claim, and
a union-bound argument.  The checked development separates these into concrete
lemmas so that the final ROM loss is not a heuristic estimate but the result of
named program invariants and probability judgments.
